%=====================================================================
\chapter{Resources}\label{app:resources}
%=====================================================================
\normalfigures

\section{The Four Books HEC Names, and Which Editions to Use}

HEC's 2017 and 2023 outlines list the same four reference books, in the same
words, with an edition for each and \textbf{no publication year for any of
them}. An edition number without a year cannot be checked against a library
catalogue or a bookshop without first finding out when it appeared --- and it
does not say whether a newer one has replaced it. The 2023 booklet adds
``(or use any other standard and latest books)'', which licenses exactly what
this appendix does: name each book at the edition to use, with its year.

\medskip
\begin{longtable}{@{}L{6.2cm}L{2.4cm}L{5.4cm}@{}}
\toprule
\rowcolor{primary!15}
\textbf{Book, as HEC lists it} & \textbf{HEC names} & \textbf{Use} \\
\midrule
\endhead
Connolly and Begg, \emph{Database Systems: A Practical Approach to Design,
Implementation, and Management} & 6th & 6th edition, Pearson, 2015 --- still
the latest \\
Garcia-Molina, Ullman and Widom, \emph{Database Systems: The Complete Book} &
2nd & 2nd edition, Pearson Prentice Hall, 2009 --- still the latest \\
Silberschatz, Korth and Sudarshan, \emph{Database System Concepts} & 6th &
\textbf{7th edition}, McGraw-Hill, 2019 --- HEC names a superseded edition \\
Ramakrishnan and Gehrke, \emph{Database Management Systems} & 3rd & 3rd
edition, McGraw-Hill, 2003 --- still the latest, and the oldest of the four \\
\bottomrule
\end{longtable}

\begin{conceptbox}[Which to read for what]
\begin{tight}
  \item \textbf{Silberschatz} is the one to own. Its chapters map most closely
        onto this course's order --- relational model, SQL, ER design,
        normalisation, indexing, transactions --- and the 7th edition covers
        the NoSQL material of \chref{project}.
  \item \textbf{Connolly and Begg} is the most patient with design: read it
        beside Part~II if \chref{ermodel} and \chref{cardinality} felt fast.
  \item \textbf{Ramakrishnan and Gehrke} is the strongest on what is inside a
        DBMS --- buffer management, B-trees, query evaluation --- and is the
        book for \chref{indexes} and \chref{transactions} if you want the
        mechanism behind the measurement.
  \item \textbf{Garcia-Molina, Ullman and Widom} is the most mathematical;
        its treatment of functional dependencies and decomposition is the one
        to read for \chref{fds} to \chref{decomposition}.
\end{tight}

\smallskip
None of the four is about SQLite. Where this handout describes SQLite's
behaviour --- affinity, \sqlkw{PRAGMA} \texttt{foreign\_keys}, file locking ---
the textbooks describe the standard, and SQLite's own documentation (below) is
the authority.
\end{conceptbox}

\section{The Tool's Own Documentation}

Every SQLite behaviour this handout measures is documented by its authors,
and the documentation is the place to settle any disagreement with a
textbook. All are at \texttt{sqlite.org}; the page names are given so they can
be found from the site's search.

\medskip
\begin{longtable}{@{}L{6.0cm}L{8.0cm}@{}}
\toprule
\rowcolor{primary!15}
\textbf{Page} & \textbf{Settles} \\
\midrule
\endhead
\emph{Datatypes In SQLite} & type affinity and storage classes ---
  \chref{ddl}'s \texttt{'hello'} in an \texttt{INTEGER} column \\
\emph{STRICT Tables} & what \sqlkw{STRICT} enforces, and since which version \\
\emph{SQLite Foreign Key Support} & why \texttt{foreign\_keys} is off by
  default, and that it is per connection --- \chref{constraints} \\
\emph{The SQLite Query Optimizer Overview} & automatic indexes, push-down,
  and the \texttt{+column} trick --- \chref{joins} to \chref{indexes} \\
\emph{EXPLAIN QUERY PLAN} & how to read \sqlkw{SCAN}, \sqlkw{SEARCH} and
  \sqlkw{TEMP B-TREE} \\
\emph{ANALYZE} & \texttt{sqlite\_stat1}, and STAT4 when it is compiled in \\
\emph{File Locking And Concurrency In SQLite Version 3} & \sqlkw{SHARED},
  \sqlkw{RESERVED}, \sqlkw{EXCLUSIVE} --- \chref{transactions} \\
\emph{Atomic Commit In SQLite} & the rollback journal and recovery after a
  crash \\
\emph{Write-Ahead Logging} & WAL mode and what it changes for readers \\
\emph{ALTER TABLE} & the four forms and the twelve-step procedure ---
  \chref{evolution} \\
\emph{JSON Functions And Operators} & \texttt{json\_object},
  \texttt{json\_each} --- \chref{project} \\
\bottomrule
\end{longtable}

\medskip
For the Python side, the \texttt{sqlite3} module's documentation at
\texttt{docs.python.org} explains \texttt{isolation\_level} --- the setting
\chref{transactions}' harness changes so that every \sqlkw{BEGIN} is one you
wrote.

\section{The Original Papers}

Three papers founded the subject. All three are short, and the first is still
the clearest statement of why a database should be relational.

\begin{tight}
  \item E.~F.~Codd, ``A Relational Model of Data for Large Shared Data
        Banks'', \emph{Communications of the ACM} 13(6), 1970 ---
        \chref{relational} and \chref{algebra}.
  \item P.~P.~Chen, ``The Entity-Relationship Model --- Toward a Unified View
        of Data'', \emph{ACM Transactions on Database Systems} 1(1), 1976 ---
        \chref{ermodel}.
  \item ANSI/X3/SPARC Study Group on Database Management Systems,
        \emph{Interim Report}, 1975 --- the three-level architecture of
        \chref{architecture}.
\end{tight}

\section{Law}

\begin{tight}
  \item \emph{Prevention of Electronic Crimes Act 2016}, Pakistan, sections 3
        and 4 --- unauthorised access to an information system and to data.
        The scope statement of Appendix~A, section A.6, and of \chref{project}
        rests on it.
\end{tight}

\section{This Handout's Own Code}

Every measured figure was produced by a file in \texttt{code/}, on the machine
described in Appendix~A, section A.5:

\medskip
\begin{longtable}{@{}L{5.2cm}L{8.8cm}@{}}
\toprule
\rowcolor{primary!15}
\textbf{File} & \textbf{Produces} \\
\midrule
\endhead
\texttt{schema/v1\_create.sql} & the university schema, from \chref{mapping} \\
\texttt{schema/v2\_constraints.sql} & the migration of \chref{evolution} \\
\texttt{schema/seed.py} & reproducible data at any scale: 4{,}000 students,
  60{,}037 enrolments by default \\
\texttt{m04\_firstdatabase.py} & \chref{firstdatabase}: time to a first row,
  bytes per page \\
\texttt{m08\_constraints.py} & \chref{constraints}: foreign keys off by default \\
\texttt{m10\_ddl.py} & \chref{ddl}: affinity against \sqlkw{STRICT} \\
\texttt{m13\_joins.py} & \chref{joins}: \sqlkw{NATURAL JOIN}, missing sections,
  automatic indexes \\
\texttt{m15\_subqueries.py} & \chref{subqueries}: anti-joins and correlation \\
\texttt{m16\_grouping.py} & \chref{grouping}: nulls, sorts, windows \\
\texttt{m18\_normalforms.py} & \chref{normalforms}: the wide table against BCNF \\
\texttt{m19\_decomposition.py} & \chref{decomposition}: lossy joins and a
  priced summary table \\
\texttt{m20\_views.py} & \chref{views}: what views and triggers cost \\
\texttt{m21\_indexes.py} & \chref{indexes}: the crossover \\
\texttt{m22\_transactions.py} & \chref{transactions}: a lost update, a
  deadlock, the price of a commit \\
\texttt{run\_all.py} & runs every harness and checks each result's
  \emph{shape} still holds \\
\texttt{sync\_labs.py} & checks that every file and command a laboratory names
  exists \\
\bottomrule
\end{longtable}
